\section*{How do Conflicts Evolve?} %previous lit
\textcolor{red}{Rework this first half to discuss civilians/protests/victimization, then build on final paragraph to write the second/latter half on frameworks that ignore interdependent nature of conflict}

More often than not, stories of armed conflict, war, or widespread criminal violence shine a spotlight on those actors who initiate injustices or wield power through collective violence against a given population. Recently, however, this narrative has begun to change. With the various uprisings across the previous 5 years--the Tunisian revolution, the Egyptian revolution, the Syrian Civil War, the Yemen revolt, the 2013 protests in Turkey, and the 2014 protests in Venezuela (just to name a few)--attention has shifted to recognize the power of the population to influence environments of extreme violence and repression. 

Compared to the number of studies focusing on economic \citep{collier04b}, political \citep{humphreys08}, and identity-based \citep{cederman10b} drivers of conflict intensity and duration, existing research has largely neglected the role of civilians in restraining or influencing the behavior of armed actors. A large body of research has explored the origins of collective action and mobilization, (Gurr 1970; Opp 1988, Tarrow 1994, Tucker 2007) and recent scholarship has assessed the effectiveness of ``maximalist campaigns" \citep{chenoweth2011} to show that countries are more likely to be democratic following nonviolent campaigns. However, as articulated by \cite{celestino2013} even the macro-level relationship between nonviolent campaigns and state-level outcomes such as democracy or regime transitions remains unclear. Importantly, \cite{celestino2013} show that nonviolent campaigns destabilize regimes, but that the trajectory of peace and democracy following such campaigns is conditional on the precise actions employed by campaign organizers. 

Building on the research agenda motivated by macro-level studies, we investigate the link between nonviolent action and violence at the local level. 

